% Realiza Barzilai–Borwein con el paso corto en una funcion modelo con
% contornos elipticos. Necesita que este definido:
%  - path(c)   : el mínimo
%  - path(x0)  : el initial guess
%  - \b        : aspect ratio de la elipse
%  - \learning : factor de aprendizaje primer paso

% Obtengo el centro
\path(c); \getlastxy
\pgfmathsetmacro{\cx}{\x}
\pgfmathsetmacro{\cy}{\y}

\fill[celeste](x0) circle(3pt) node[above right] {$x_0$};
               
% Posicion
\path(x0); \getlastxy

% Calculo el gradiente
\pgfmathsetmacro{\dx}{2.*(\x-\cx)} % 2(x-x0)/a^2
\pgfmathsetmacro{\dy}{2.*(\y-\cy)/(\b*\b)} % 2(y-y0)/b^2

% Primer paso con leraning constante
\pgfmathsetmacro{\norm}{sqrt(\dx*\dx+\dy*\dy)}
\pgfmathsetmacro{\learning}{\learning/\norm}
               
% Move
\pgfmathsetmacro{\vx}{-\learning*\dx}
\pgfmathsetmacro{\vy}{-\learning*\dy}
\draw[->,shorten >=5, thick, amarillo] (x0) --  ++(\vx,\vy);
\path(x0) ++(\vx,\vy) coordinate(x1);
\fill[celeste](x1) circle(3pt) node[above right] {$x1$};
 
% Save last position as the new x0
\path(x1) coordinate(x0);       
        
% Save last gradient
\pgfmathsetmacro{\dxx}{\dx}
\pgfmathsetmacro{\dyy}{\dy}

\foreach \m in {2,...,25} {

  % Posicion
  \path(x0); \getlastxy

  % Gradiente
  \pgfmathsetmacro{\dx}{2.*(\x-\cx)} % 2(x-x0)/a^2
  \pgfmathsetmacro{\dy}{2.*(\y-\cy)/(\b*\b)} % 2(y-y0)/b^2
 
  % Compute learning rate from previous iteration information
  \pgfmathsetmacro{\ddx}{\dx-\dxx}
  \pgfmathsetmacro{\ddy}{\dy-\dyy}
  % \pgfmathsetmacro{\learning}{(\ddx*\vx+\ddy*\vy)/(\ddx*\ddx+\ddy*\ddy)} % short
  \pgfmathsetmacro{\learning}{(\vx*\vx+\vy*\vy)/(\ddx*\vx+\ddy*\vy)} % long
                            
  % Move
  \pgfmathsetmacro{\vx}{-\learning*\dx}
  \pgfmathsetmacro{\vy}{-\learning*\dy}
  \path(x0) ++(\vx,\vy) coordinate(x\m);
  \ifthenelse{\m<5}{
    \draw[->,shorten >=5, thick, amarillo] (x0) --  ++(\vx,\vy);
    \fill[celeste](x\m) circle(3pt) node[above] {$x_\m$};
  }{
    \fill[celeste](x\m) circle(3pt);
  }

  % Save last position as the new x0
  \path(x\m) coordinate(x0);
 
  % Save last gradient
  \pgfmathsetmacro{\dxx}{\dx}
  \pgfmathsetmacro{\dyy}{\dy}
               
}

 
% Contornos de la funcion  
\foreach \i in {1,2,...,11} {
\pgfmathsetmacro\lab{\i+0.1}
\scriptsize\draw[rosa,thick] (c) ellipse  [x radius=\i*1cm, y radius=\i*\b*1cm] node[xshift=\lab cm]{\i};
}
     
